%% RESUMO NA LINGUA VERNACULA (elemento obrigatorio -- NBR 6028)
%%
%% Regras:
%%   - paragrafo unico, sem recuo, justificado;
%%   - Times New Roman 12, espacamento simples;
%%   - entre 150 e 500 palavras;
%%   - verbo na voz ativa e na terceira pessoa do singular;
%%   - as palavras-chave sao definidas em config/dados.tex.
%%
%% Um bom resumo informa: objetivo geral; principais conceitos e autores da
%% fundamentacao teorica; etapas e acoes da execucao; populacao envolvida e
%% periodo; principais resultados alcancados; conclusoes.

\begin{resumo}
Esta pesquisa investiga a aplicação de modelos estatísticos espaço-temporais, técnicas de machine learning e abordagens de big data sobre os dados do Sistema Nacional de Pesquisa de Custos e Índices da Construção Civil (SINAPI), com o objetivo de apoiar a tomada de decisão em políticas públicas de infraestrutura. O trabalho parte do reconhecimento de que o SINAPI constitui a principal referência nacional de custos da construção civil e apresenta potencial ainda pouco explorado para análises preditivas. A pesquisa busca responder à importância da previsão dos dados do SINAPI e avaliar se a utilização de valores de referência da região de São Paulo para localidades sem observações próprias pode gerar sobrepreço em contratações públicas. A fundamentação teórica aborda conceitos de ciência de dados, inteligência artificial, machine learning, análise espaço-temporal e gestão baseada em evidências para políticas públicas. A proposta contempla a construção e a validação de modelos capazes de identificar tendências, estimar custos futuros e apoiar o planejamento orçamentário de investimentos em habitação, saneamento e infraestrutura. Também são analisados os benefícios da antecipação de riscos relacionados à insuficiência orçamentária, à paralisação de obras e ao desequilíbrio econômico-financeiro de contratos. Os dados utilizados são provenientes da base pública do SINAPI, mantida pela CAIXA Econômica Federal em parceria com o Instituto Brasileiro de Geografia e Estatística (IBGE), abrangendo séries históricas de custos e insumos da construção civil. Como resultado esperado, pretende-se demonstrar que a aplicação de técnicas avançadas de análise de dados pode transformar o SINAPI de uma ferramenta predominantemente consultiva em uma plataforma de inteligência capaz de apoiar decisões estratégicas, aumentar a eficiência da alocação de recursos públicos e fortalecer a formulação de políticas públicas baseadas em evidências.
\end{resumo}
